\section{Related Work}

Table~\ref{tab:overview} (Appendix~\ref{app:relwork-table}) surveys prior work across three related areas; we explain the key distinctions below.

\subsection{Agentic Data and Tool-Calling Benchmarks}
The dominant paradigm for evaluating tool-calling agents uses either environment-based execution feedback~\cite{zhou2024webarena,drouin2024workarena,trivedi2024appworld} or deterministic scoring against annotated trajectories~\cite{deng2023mind2web,xu2025agenttrek,patil2025bfcl}. BFCL~\cite{patil2025bfcl} is the closest structural ancestor to our data format: it isolates single-turn function-calling using typed JSON schemas and AST-based ground-truth comparison. We adopt the same schema but extend it to multi-step DAG-structured workflows at three controlled difficulty tiers. The most structurally similar work is FuncBenchGen~\cite{funcbenchgen2025}, which frames multi-step calling as DAG traversal with controllable complexity. The key distinction is purpose: FuncBenchGen was designed to train and evaluate \emph{generator} agents, whereas AgentJudgeBench is designed to measure the reliability of \emph{judge} models applied to those generators' outputs - a different and complementary objective that requires a paired GT/no-GT evaluation protocol and per-metric decomposition that FuncBenchGen does not provide. TaskBench~\cite{shen2024taskbench} shares the tool-dependency graph framing but focuses on decomposition quality rather than judge alignment. Synthetic data pipelines~\cite{wang2022selfinstruct,xu2024magpie,cui2024ultrafeedback} inform our generation approach, but target model training rather than evaluation benchmarking.

\subsection{LLM-as-Judge}
Zheng et al.~\cite{zheng2023judging} established the LLM-as-judge paradigm on MT-Bench, demonstrating GPT-4's strong agreement with human preferences on open-ended dialogue while identifying positional preference, verbosity sensitivity, and self-enhancement as systematic biases. Subsequent work either extends the evaluation surface - JudgeBench~\cite{tan2024judgebench} and Arena-Hard-Auto~\cite{li2024arenahard} move to hard, verifiable response pairs - or builds dedicated judge models~\cite{wang2024pandalm,li2024autoj}. The jury-of-judges paradigm~\cite{verga2024jury} shows that ensembling reduces individual judge biases. Crucially, all of this work targets text-centric tasks where judge quality reduces to preference alignment over fluent, semantically coherent outputs. AgentJudgeBench occupies a different regime: correctness is structural, the verdict space is $\{0, 0.5, 1\}$ rather than a preference ranking, and failure modes are multi-dimensional and interdependent.

\subsection{LLM Judges for Tool-Calling}
LLM judges have been deployed for tool-calling evaluation in several recent benchmarks, but always as an implementation detail rather than the subject of study. ToolEval~\cite{qin2023toolllm} uses ChatGPT to compute pass rates over API trajectories, reporting 87.1\% human agreement in aggregate; StableToolBench~\cite{guo2024stabletoolbench} and MCP-AgentBench~\cite{mcpagentbench2025} adopt similar approaches with updated judge models. GeoBenchX~\cite{krechetova2025geobenchx} assembles a three-judge panel achieving 88-96\% agreement. Agent-as-a-Judge~\cite{zhuge2024agentjudge} and Auto-Eval Judge~\cite{bhonsle2025autoeval} propose more structured evaluation frameworks. In each case, judge reliability is reported as a single aggregate figure on a small validation sample, with no variation of task complexity, difficulty tier, or ground-truth availability. ToolSandbox~\cite{lu2024toolsandbox} is notable for questioning LLM judge reliability directly and replacing judges with programmatic evaluation - but it replaces the judge rather than characterising its failure modes. AgentJudgeBench takes the complementary stance: keep the judge paradigm and measure it systematically, decomposing reliability by metric, topology, difficulty, and condition across $321{,}648$ completed evaluations (Appendix~\ref{app:eval-stats}). A detailed feature comparison across all seven related systems is in Appendix~\ref{app:judge-comparison}.
